\documentclass[10pt,a4paper]{amsart}
\usepackage[margin=19mm]{geometry}
\usepackage{amsmath,amssymb,mathtools}
\usepackage[scr=rsfs,cal=euler]{mathalfa}
\usepackage{xfrac}
\usepackage[unicode,hidelinks,bookmarks=false]{hyperref}
\newcommand{\ie}{i.e.\ }
\newcommand{\cf}{cf.\ }
\newcommand{\resp}{respectively\ }
\newcommand{\Subsection}{\subsection}
\providecommand{\nobreakdash}{\nobreak\hbox{-}}
\setlength{\emergencystretch}{2em}
\multlinegap0pt
\usepackage{bm}
\usepackage{bbm}
\usepackage{dsfont}
\usepackage{xspace}
\usepackage{cancel}
\usepackage{mathtools}
\usepackage{amsthm}
\makeatletter
\@ifundefined{proposition}{\newtheorem{proposition}{Proposition}}{}
\makeatother
\title[Hamiltonian formulation of the quasineutral Vlasov-Poisson s]{Hamiltonian formulation of the quasineutral Vlasov-Poisson system}
\author{J. W. Burby, D. A. Kaltsas, P. J. Morrison, E. Tassi,, G. N. Throumoulopoulos}
\date{Source: arXiv:2506.21415v2 (CC BY 4.0)}
\begin{document}
\maketitle

\noindent\emph{Source and licence.} Hamiltonian formulation of the quasineutral Vlasov-Poisson system. Version arXiv:2506.21415v2 (CC BY 4.0), distributed under the Creative Commons Attribution 4.0 International licence, as recorded in the arXiv licence metadata for this exact version and shown on the abstract page https://arxiv.org/abs/2506.21415v2. Full rights notice: licenses/arxiv-2506.21415-CC-BY-4.0.txt.\par\smallskip

\noindent\textit{Editorial scope.} Counted unit: the Proposition whose source label is \texttt{PD\_proposition} in the article (source0.tex lines 326--332, source label \texttt{PD\_proposition}), delivered together with the article's own complete original proof of it (source0.tex lines 338--402, one \texttt{proof} environment of 65 lines). No mathematics is written, repaired or re-derived: statement and proof are the article's own text line for line. The proof is detached from the statement in the source: the article states the result at lines 326--332 and prints its proof at lines 338--402, and the proof environment's own header names the statement, so the association is the article's own. Required context retained uncounted: VP\_cbracket (lines 266--276); Pdot\_G (lines 286--295); ndot\_G (lines 286--295); varrhodot\_G (lines 286--295). Every definition, display and label of those lines is retained unchanged. Omitted from this package and not redistributed: 1--265, 277--285, 296--325, 333--337, 403--563. Nothing inside a retained excerpt is omitted or altered: every retained body is delivered exactly as the article's lines are written, so no omission receipt of any kind accompanies this package. Citations: the retained excerpts carry no citation command at all, so no bibliography entry of the article is transcribed and no reference is redistributed. Reference adaptation: none was needed and none was made; every reference inside the delivered unit resolves inside it, so no unresolved reference or citation is produced. Printed slips kept literal: no printed word, symbol, hypothesis or number of the counted statement or of its proof was altered. Layout and class: the article's own class is \texttt{revtex4-1} with packages including amsthm, graphicx, dcolumn, bm, inputenc, fontenc, mathptmx, etoolbox, babel; the wrapper is the lane's pinned \texttt{amsart} editorial wrapper at 10pt on A4, which restates the theorem-like environments the retained text needs and adds the article's own macro definitions that the retained text needs (none), with extra wrapper packages (bm, bbm, dsfont, xspace, cancel, mathtools). The delivered numbering is the unit's own. Assets: no retained span contains a figure environment, a graphics-inclusion command, a PDF-page inclusion, a table or a TikZ picture; no image file is copied into this package, the package declares no asset dependency and no image file is redistributed with it. Licence: version arXiv:2506.21415v2 of this article is distributed under the Creative Commons Attribution 4.0 International licence, as recorded in the arXiv licence metadata for this exact version and shown on the abstract page https://arxiv.org/abs/2506.21415v2. A full rights notice is delivered at licenses/arxiv-2506.21415-CC-BY-4.0.txt. Characters: every retained excerpt is pure ASCII, so the delivered engine, XeLaTeX with its default Latin Modern text font, reports no missing character.
\begin{align}
&\{F,G\}_{E} =-\epsilon\int\bigg(\frac{\delta F}{\delta\bm{P}}\cdot \bigg[\partial_{\bm{q}}\frac{\delta G}{\delta n} - \partial_{\bm{q}}\left\langle\frac{1}{n}\frac{\delta G}{\delta\varrho}\right\rangle + \left\langle\partial_{\bm{q}}\frac{1}{n}\frac{\delta G}{\delta\varrho}\right\rangle\bigg]-\frac{\delta G}{\delta\bm{P}}\cdot \bigg[\partial_{\bm{q}}\frac{\delta F}{\delta n} - \partial_{\bm{q}}\left\langle\frac{1}{n}\frac{\delta F}{\delta\varrho}\right\rangle + \left\langle\partial_{\bm{q}}\frac{1}{n}\frac{\delta F}{\delta\varrho}\right\rangle\bigg]\bigg)\,n\,d\bm{q}\nonumber\\
%
&-\epsilon\int\bigg[\frac{\delta F}{\delta\bm{P}},\frac{\delta G}{\delta\bm{P}}\bigg]\cdot\bm{P}\,d\bm{q} +\int B\,\frac{\delta F}{\delta\bm{P}}\cdot \mathbb{J}\frac{\delta G}{\delta\bm{P}} \,n\,d\bm{q}\nonumber\\
%%%
&+\epsilon\int \bigg(\partial_{\bm{q}}\frac{1}{n}\frac{\delta F}{\delta \varrho}\cdot \partial_{\bm{\xi}}\frac{1}{n}\frac{\delta G}{\delta\varrho} -\partial_{\bm{q}}\frac{1}{n}\frac{\delta G}{\delta \varrho}\cdot \partial_{\bm{\xi}}\frac{1}{n}\frac{\delta F}{\delta\varrho} \bigg)\,n\,\varrho\,d\bm{\xi}\,d\bm{q} + \int (B-\epsilon\,\Omega)\bigg(\partial_{\bm{\xi}}\frac{1}{n}\frac{\delta F}{\delta\varrho}\bigg)\cdot\mathbb{J}\bigg(\partial_{\bm{\xi}}\frac{1}{n}\frac{\delta G}{\delta\varrho}\bigg)\,n\,\varrho\,d\bm{\xi}\,d\bm{q}\nonumber\\
%%%
&-\epsilon\int \bigg(\left\langle\partial_{\bm{q}}\frac{1}{n}\frac{\delta F}{\delta\varrho}\right\rangle\cdot\left\langle\partial_{\bm{\xi}}\frac{1}{n}\frac{\delta G}{\delta\varrho}\right\rangle-\left\langle\partial_{\bm{q}}\frac{1}{n}\frac{\delta G}{\delta\varrho}\right\rangle\cdot\left\langle\partial_{\bm{\xi}}\frac{1}{n}\frac{\delta F}{\delta\varrho}\right\rangle\bigg)\,n\,d\bm{q}-\int (B-\epsilon\,\Omega)\, \left\langle\partial_{\bm{\xi}}\frac{1}{n}\frac{\delta F}{\delta\varrho}\right\rangle\cdot \mathbb{J} \left\langle\partial_{\bm{\xi}}\frac{1}{n}\frac{\delta G}{\delta\varrho}\right\rangle\,n\,d\bm{q}\nonumber\\
%%%
&- \epsilon \int \bigg(\left[\frac{\delta F}{\delta\bm{P}} - \left\langle\partial_{\bm{\xi}}\frac{1}{n}\frac{\delta F}{\delta\varrho}\right\rangle\right]\cdot\partial_{\bm{q}}\cdot\left[n\left\langle\partial_{\bm{\xi}}\left(\frac{1}{n}\frac{\delta G}{\delta\varrho}\right)\bm{\xi}\right\rangle\right] -\left[\frac{\delta G}{\delta\bm{P}} - \left\langle\partial_{\bm{\xi}}\frac{1}{n}\frac{\delta G}{\delta\varrho}\right\rangle\right]\cdot\partial_{\bm{q}}\cdot\left[n\left\langle\partial_{\bm{\xi}}\left(\frac{1}{n}\frac{\delta F}{\delta\varrho}\right)\bm{\xi}\right\rangle\right] \bigg)\,d\bm{q}.\label{VP_cbracket}
\end{align}

\begin{align}
&\partial_tn +\partial_{\bm{q}}\cdot\left(n\frac{\delta G}{\delta\bm{P}}\right) = 0 \label{ndot_G}\\
%%%
&\partial_t\bm{P} + \epsilon\,n\,\partial_{\bm{q}}\left(\frac{\delta G}{\delta\bm{P}}\cdot \frac{\bm{P}}{n}\right) + \epsilon\,\partial_{\bm{q}}\cdot\left(n\,\frac{\delta G}{\delta\bm{P}}\right)\frac{\bm{P}}{n} + \epsilon\,n\,\Omega\mathbb{J}\frac{\delta G}{\delta \bm{P}} =-\epsilon\,\partial_{\bm{q}}\cdot\left(n\left\langle\partial_{\bm{\xi}}\left(\frac{1}{n}\frac{\delta G}{\delta\varrho}\right)\bm{\xi} \right\rangle\right) \nonumber\\
&\hspace*{21em}- \epsilon\,n\,\bigg(\partial_{\bm{q}}\frac{\delta G}{\delta n} - \partial_{\bm{q}}\left\langle \frac{1}{n}\frac{\delta G}{\delta\varrho}\right\rangle + \left\langle\partial_{\bm{q}}\frac{1}{n}\frac{\delta G}{\delta\varrho} \right\rangle\bigg) + n\,B\,\mathbb{J}\frac{\delta G}{\delta\bm{P}} \label{Pdot_G}\\
%%%
&\partial_t(n\varrho) + \partial_{\bm{q}}\cdot\bigg(\bigg[\epsilon\,\frac{\delta G}{\delta\bm{P}} +\epsilon\partial_{\bm{\xi}}\frac{1}{n}\frac{\delta G}{\delta\varrho}- \epsilon\left\langle\partial_{\bm{\xi}}\frac{1}{n}\frac{\delta G}{\delta\varrho}\right\rangle\bigg]\,n\varrho\bigg)\nonumber\\
&+ \partial_{\bm{\xi}}\cdot\bigg(\bigg[-\epsilon\,\partial_{\bm{q}}\left(\bm{\xi}\cdot\frac{\delta G}{\delta\bm{P}}\right)+\epsilon\,\partial_{\bm{q}}\left(\bm{\xi}\cdot \left\langle\partial_{\bm{\xi}}\frac{1}{n}\frac{\delta G}{\delta\varrho}\right\rangle\right) + \epsilon\,n^{-1}\,\partial_{\bm{q}}\cdot\left(n\left\langle\partial_{\bm{\xi}}\left(\frac{1}{n}\frac{\delta G}{\delta\varrho}\right)\bm{\xi}\right\rangle\right)\nonumber\\
&\qquad\qquad-\epsilon\,\partial_{\bm{q}}\left(\frac{1}{n}\frac{\delta G}{\delta\varrho}\right) + \epsilon\,\left\langle\partial_{\bm{q}}\left(\frac{1}{n}\frac{\delta G}{\delta\varrho}\right)\right\rangle + (B - \epsilon\Omega)\mathbb{J}\bigg(\partial_{\bm{\xi}}\frac{1}{n}\frac{\delta G}{\delta \varrho} - \left\langle \partial_{\bm{\xi}}\frac{1}{n}\frac{\delta G}{\delta \varrho}\right\rangle\bigg)\bigg]\,n\varrho\bigg) = 0.\label{varrhodot_G}
\end{align}

\begin{proposition}\label{PD_proposition}
Let $\mathfrak{s}^*$ denote the space of tuples $(n,\bm{P},\varrho)$, where $n:\mathbb{T}^2\rightarrow\mathbb{R}$, $\bm{P}:\mathbb{T}^2\rightarrow\mathbb{R}$, and $\varrho:\mathbb{T}^2\times\mathbb{R}^2\rightarrow\mathbb{R}$. The submanifold $\Sigma\subset\mathfrak{s}^*$ defined by 
\begin{align*}
\Sigma = \{(n,\bm{P},\varrho)\in\mathfrak{s}^*\mid \partial_{\bm{q}}n = 0,\quad \partial_{\bm{q}}\cdot\bm{P} = 0\}
\end{align*}
is a Poisson-Dirac submanifold when $\mathfrak{s}^*$ is equipped with the Poisson bracket $\{\cdot,\cdot\}_E$ from Eq.\,\eqref{VP_cbracket}. In particular, there is a Poisson bracket $\{\cdot,\cdot\}_\Sigma$ on $\Sigma$ induced by $\{\cdot,\cdot\}_E$ given explicitly in Eq.\,\eqref{QNVP_cbracket}.
\end{proposition}

\begin{proof}[Proof of Proposition \ref{PD_proposition}]
To establish that $\Sigma$ is Poisson-Dirac, we must show (A) that it satisfies the Poisson-Dirac condition $T_\sigma \Sigma \cap \widehat{\pi}_E(T_\sigma\Sigma^\circ) = \{0\}$, and (B) that the induced bracket guaranteed by (A) varies smoothly along $\Sigma$. Here $\widehat{\pi}_E$ denotes the bundle map $T^*\mathfrak{s}^*\rightarrow T\mathfrak{s}^*$ associated with $\{\cdot,\cdot\}_E$ and $T_\sigma\Sigma^\circ\subset T^*_\sigma\mathfrak{s}^*$ denotes the annihilator of $T_\sigma\Sigma\subset T_\sigma\mathfrak{s}^*$. 

(A) Let $\sigma = (n_0,\bm{\pi},\varrho)\in\Sigma$ be a general point in $\Sigma$. Here, $n_0$ is a real constant and $\bm{\pi}$ is a divergence-free vector field on $\mathbb{T}^2$. The tangent space to $\Sigma$ at $\sigma$ is given by 
\begin{align*}
T_\sigma\Sigma = \{(\delta n_0,\delta\bm{\pi},\delta\varrho)\mid \delta n_0\in \mathbb{R},\quad \delta\bm{\pi}:\mathbb{T}^2\rightarrow\mathbb{R}^2,\quad \delta\varrho:\mathbb{T}^2\times \mathbb{R}^2\rightarrow\mathbb{R}, \quad \partial_{\bm{q}}\cdot\delta\bm{\pi} = 0\}.
\end{align*}
The annihilator of $T_\sigma\Sigma$ in $T^*_\sigma\mathfrak{s}^*\ni (\delta n^*,\delta\bm{P}^*,\delta\varrho^*)$ is therefore
\begin{align*}
T_\sigma\Sigma^\circ = \{(\delta n^*,\delta\bm{P}^*,\delta\varrho^*)\in T^*_\sigma\mathfrak{s}^*\mid \exists\delta\Phi^*:\mathbb{T}^2\rightarrow\mathbb{R},\quad\int\delta\Phi^*\,d\bm{q} = 0,\quad \int \delta n^*\,d\bm{q} = 0,\quad \delta\bm{P}^* =\partial_{\bm{q}}\delta\Phi^*,\quad \delta\varrho^* = 0\}.
\end{align*}
Suppose that the tangent vector $\delta z =(\delta n,\delta\bm{P},\delta\varrho) \in T_\sigma\mathfrak{s}^*$ is contained in both $T_\sigma\Sigma$ and $\widehat{\pi}_E(T_\sigma\Sigma^\circ)$. Since $\delta z\in \widehat{\pi}_E(T_\sigma\Sigma^\circ)$ there must be $\delta\alpha=(\delta n^*,\delta\bm{P}^*,\delta\varrho^*) = (\delta n^*,\partial_{\bm{q}}\delta\Phi^*,0)\in T_\sigma\Sigma^\circ$ such that $\delta z = \widehat{\pi}_E(\delta \alpha)$. Notice that $\widehat{\pi}_E(\delta \alpha)$ can be computed explicitly using Eqs.\,\eqref{ndot_G}-\eqref{varrhodot_G} by making the substitutions $\delta G/\delta n \rightarrow \delta n^*$, $\delta G/\delta\bm{P} \rightarrow \delta\bm{P}^*$, and $\delta G/\delta\varrho \rightarrow \delta\varrho^*$. By Eq.\,\eqref{ndot_G} this implies $\delta n = -\partial_{\bm{q}}\cdot(n_0\partial_{\bm{q}}\delta\Phi^*) = - n_0\,\Delta\delta\Phi^*$. On the other hand, $\delta z\in T_\sigma \Sigma$ implies that $\delta n = \delta n_0$ must be constant. Integrating $\delta n_0 = - n_0\,\Delta \delta\Phi^*$ over $\mathbb{T}^2$ implies that constant must vanish, $\delta n_0 = 0$. Thus, $\Delta \delta\Phi^* = 0$, which implies $\delta\Phi^* = 0$ because $\delta\Phi^*$ has zero mean. It follows that $\delta n = -n_0\Delta\delta\Phi^* = 0$. By Eq.\,\eqref{Pdot_G}, $\delta z = \widehat{\pi}_E(\delta\alpha)$ also implies $\delta\bm{P} = -\epsilon\,n_0\partial_{\bm{q}}\delta n^*$. But $\delta z\in T_\sigma\Sigma$ implies $\delta\bm{P} = \delta\bm{\pi}$ is divergence-free. Therefore $\Delta \delta n^* = 0$, which requires $\delta n^* = 0$ because $\delta n^*$ has zero mean. It follows that $\delta\bm{P} = - \epsilon\,n_0\,\partial_{\bm{q}}\delta n^* = 0$. Finally, in light of $\delta n^* = \delta\Phi^* = 0$ and $\delta z = \widehat{\pi}_E(\delta\alpha)$, Eq.\,\eqref{varrhodot_G} implies $\delta\varrho = 0$. We have shown $\delta z \in T_\sigma\Sigma\cap \widehat{\pi}_E(T_\sigma\Sigma^\circ)$ implies $\delta z = 0$, which establishes the claim $T_\sigma\Sigma\cap \widehat{\pi}_E(T_\sigma\Sigma^\circ) = \{0\}$ for all $\sigma\in\Sigma$.

(B) Part (A) of the proof guarantees that there is a bilinear bracket operation $\{\cdot,\cdot\}_\Sigma$ between functionals $F,G:\Sigma\rightarrow\mathbb{R}$. To complete the proof we must demonstrate that $\{F,G\}_\Sigma$ is a smooth function on $\Sigma$ when $F,G$ are both smooth. For this demonstration we will find and analyze an explicit formula for $\{F,G\}_\Sigma$. The value of the bracket $\{F,G\}_\Sigma$ at $\sigma=(n_0,\bm{\pi},\varrho)\in \Sigma$ is given by $\{F,G\}_\Sigma(\sigma) = \widetilde{(dF_\sigma)}\cdot \pi_{E}\cdot \widetilde{(dG_\sigma)}$, where $\pi_E$ denotes the Poisson tensor on $\mathfrak{s}^*$ associated with $\{\cdot,\cdot\}_E$. Here $dF_\sigma,dG_\sigma\in T^*_\sigma\Sigma$ denote the differentials of $F,G$ at the point $\sigma$ and $\widetilde{(dF_\sigma)},\widetilde{(dG_\sigma)}\in T^*_\sigma\mathfrak{s}^*$ are any covectors in $T^*_\sigma\mathfrak{s}^*$ such that
\begin{align}
&\widetilde{(dF_\sigma)}\mid T_\sigma\Sigma = dF_\sigma,\quad \widetilde{(dG_\sigma)}\mid T_\sigma\Sigma = dG_\sigma\label{extension_property}\\
&\widetilde{(dF_\sigma)}\mid \widehat{\pi}_E(T_\sigma\Sigma^\circ) = 0,\quad \widetilde{(dG_\sigma)}\mid \widehat{\pi}_E(T_\sigma\Sigma^\circ) = 0.\label{transverse_property}
\end{align}
When applied to $\delta \sigma = (\delta n_0,\delta\bm{\pi},\delta\varrho)\in T_\sigma\Sigma$ the value of $dG_\sigma$ is
\begin{align*}
dG_\sigma(\delta \sigma) = \int \frac{\delta G}{\delta n_0}\,\delta n_0\,d\bm{q} + \int \frac{\delta G}{\delta \bm{\pi}}\cdot\delta\bm{\pi}\,d\bm{q} + \int\frac{\delta G}{\delta\varrho}\,\delta\varrho\,d\bm{\xi}\,d\bm{q},
\end{align*}
where $\delta G/\delta n_0$ is required to be constant and $\delta G/\delta\bm{\pi}$ is required to be divergence-free.
Suppose that the value of $\widetilde{dG_\sigma}$ when applied to $\delta z = (\delta n,\delta \bm{P},\delta\varrho)\in T_\sigma\mathfrak{s}^*$ is given by
\begin{align*}
\widetilde{dG_\sigma}(\delta z) = \int \widetilde{\frac{\delta G}{\delta n_0}}\,\delta n\,d\bm{q} + \int \widetilde{\frac{\delta G}{\delta\bm{\pi}}}\cdot \delta\bm{P}\,d\bm{q} + \int \widetilde{\frac{\delta G}{\delta\varrho}}\,\delta\varrho\,d\bm{\xi}\,d\bm{q},
\end{align*}
where $\widetilde{\frac{\delta G}{\delta n_0}}$, $\widetilde{\frac{\delta G}{\delta\bm{\pi}}}$, and $\widetilde{\frac{\delta G}{\delta\varrho}}$ are unknown coefficients. The condition \eqref{extension_property} implies
\begin{align*}
\widetilde{\frac{\delta G}{\delta n_0}}  = \frac{\delta G}{\delta n_0} + \psi_G,\quad \widetilde{\frac{\delta G}{\delta\bm{\pi}}}  = \frac{\delta G}{\delta\bm{\pi}} + \partial_{\bm{q}}\Phi_G,\quad \widetilde{\frac{\delta G}{\delta\varrho}}  = \frac{\delta G}{\delta \varrho},
\end{align*}
where $\psi_G$ and $\Phi_G$ are both functions on $\mathbb{T}^2$ with zero mean, not determined by condition \eqref{extension_property}. In order for $\widetilde{dG_\sigma}$ to satisfy property \eqref{transverse_property}, for each $\alpha = (\delta n^*,\partial_{\bm{q}}\delta\Phi^*,0)\in T_\sigma\Sigma^0$ the tangent vector $\delta z = \widehat{\pi}_E(\alpha)\in T_\sigma\mathfrak{s}^*$ must annihilate $\widetilde{dG_\sigma}$. An explicit expression for $\delta z = (\delta n,\delta\bm{P},\delta\varrho)$ follows from Eqs.\,\eqref{ndot_G}-\eqref{varrhodot_G} using the substitutions $\delta G/\delta n\rightarrow \delta n^*$, $\delta G/\delta \bm{P}\rightarrow \partial_{\bm{q}}\delta \Phi^*$, and $\delta G/\delta\varrho\rightarrow 0$. The result is
\begin{align}
&\delta n +\epsilon\partial_{\bm{q}}\cdot\left(n_0\partial_{\bm{q}}\delta\Phi^*\right) = 0 \label{delta_n_biv}\\
%%%
&\delta\bm{P} + \epsilon\,\partial_{\bm{q}}\left( {\bm{\pi}}\cdot \partial_{\bm{q}}\delta\Phi^*\right) + \epsilon\,\partial_{\bm{q}}\cdot\left(\partial_{\bm{q}}\delta\Phi^*\right){\bm{\pi}} + \epsilon\,n_0\,\Omega\mathbb{J}\partial_{\bm{q}}\delta\Phi^* = - \epsilon\,n_0\,\partial_{\bm{q}}\delta n^* + n_0\,B\,\mathbb{J}\partial_{\bm{q}}\delta\Phi^* \label{delta_P_biv}\\
%%%
&\delta\varrho + \bigg[\epsilon\,\partial_{\bm{q}}\delta\Phi^* \bigg]\cdot\partial_{\bm{q}}\varrho- \partial_{\bm{\xi}}\cdot\bigg(\bigg[\epsilon\,\partial_{\bm{q}}\left(\bm{\xi}\cdot\partial_{\bm{q}}\delta\Phi^*\right)\bigg]\varrho\bigg) = 0.\label{delta_varrho_biv}
\end{align}
The condition $\widetilde{dG_\sigma}(\delta z) = 0$ therefore implies
\begin{align*}
&\Delta \Phi_G = 0\\
&\epsilon\,\Delta \psi_G = \partial_{\bm{q}}\cdot\bigg([B -\epsilon\,\Omega]\mathbb{J}\frac{\delta G}{\delta\bm{\pi}}\bigg) - \epsilon\,\Delta\left(\frac{\bm{\pi}}{n_0}\cdot \frac{\delta G}{\delta\bm{\pi}}\right) - \epsilon\,\partial_{\bm{q}}\partial_{\bm{q}}:\left\langle\left(\partial_{\bm{\xi}}\frac{1}{n_0}\frac{\delta G}{\delta \varrho}\right)\bm{\xi}\right\rangle + \epsilon\,\partial_{\bm{q}}\cdot\left\langle \frac{1}{n_0}\frac{\delta G}{\delta \varrho}\partial_{\bm{q}}\ln\varrho\right\rangle,
\end{align*}
which uniquely determines $\Phi_G$ and $\psi_G$. Note in particular that $\Phi_G = 0$. The bracket $\{F,G\}_\Sigma = \widetilde{dF_\sigma}\cdot\pi_E\cdot\widetilde{dG_\sigma}$ may now be computed explicitly by making the following substitutions in Eq.\,\eqref{VP_cbracket}:
\begin{align*}
\frac{\delta G}{\delta n}&\rightarrow \frac{\delta G}{\delta n_0} -\frac{\bm{\pi}}{n_0}\cdot\frac{\delta G}{\delta\bm{\pi}} + \mathcal{G}\bigg(\partial_{\bm{q}}\cdot\bigg[\frac{1}{\epsilon}(B-\epsilon\,\Omega)\mathbb{J}\frac{\delta G}{\delta\bm{\pi}} - \partial_{\bm{q}}\cdot\left(\left\langle\partial_{\bm{\xi}}\left(\frac{1}{n_0}\frac{\delta G}{\delta\varrho}\right)\bm{\xi}\right\rangle\right) + \partial_{\bm{q}}\left\langle\frac{1}{n_0}\frac{\delta G}{\delta \varrho}\right\rangle - \left\langle\partial_{\bm{q}}\frac{1}{n_0}\frac{\delta G}{\delta\varrho}\right\rangle\bigg]\bigg)+\text{const.}\\
\frac{\delta G}{\delta\bm{P}}&\rightarrow \frac{\delta G}{\delta\bm{\pi}}\\
\frac{\delta G}{\delta\varrho} & \rightarrow\frac{\delta G}{\delta\varrho}.
\end{align*}
The result is
\begin{align}
&\{F,G\}_{\Sigma} = \int (B-\epsilon\,\Omega)\,\frac{\delta F}{\delta\bm{\pi}}\cdot \mathbb{J}\frac{\delta G}{\delta\bm{\pi}} \,n_0\,d\bm{q} \nonumber\\
%
&-\epsilon\int\bigg(\frac{\delta F}{\delta\bm{\pi}}\cdot \bigg[    \left\langle\partial_{\bm{q}}\frac{\delta G}{\delta\varrho}\right\rangle\bigg]-\frac{\delta G}{\delta\bm{\pi}}\cdot \bigg[  \left\langle\partial_{\bm{q}}\frac{\delta F}{\delta\varrho}\right\rangle\bigg]\bigg)\,d\bm{q}\nonumber\\
%%%
&- \epsilon \int \bigg(\left[\frac{\delta F}{\delta\bm{\pi}} - \left\langle\partial_{\bm{\xi}}\frac{1}{n_0}\frac{\delta F}{\delta\varrho}\right\rangle\right]\cdot\partial_{\bm{q}}\cdot\left[\left\langle\partial_{\bm{\xi}}\left(\frac{\delta G}{\delta\varrho}\right)\bm{\xi}\right\rangle\right] -\left[\frac{\delta G}{\delta\bm{\pi}} - \left\langle\partial_{\bm{\xi}}\frac{1}{n_0}\frac{\delta G}{\delta\varrho}\right\rangle\right]\cdot\partial_{\bm{q}}\cdot\left[\left\langle\partial_{\bm{\xi}}\left(\frac{\delta F}{\delta\varrho}\right)\bm{\xi}\right\rangle\right] \bigg)\,d\bm{q}\nonumber\\
%%%
&+\epsilon\int \bigg(\partial_{\bm{q}}\frac{1}{n_0}\frac{\delta F}{\delta \varrho}\cdot \partial_{\bm{\xi}}\frac{1}{n_0}\frac{\delta G}{\delta\varrho} -\partial_{\bm{q}}\frac{1}{n_0}\frac{\delta G}{\delta \varrho}\cdot \partial_{\bm{\xi}}\frac{1}{n_0}\frac{\delta F}{\delta\varrho} \bigg)\,n_0\,\varrho\,d\bm{\xi}\,d\bm{q} + \int (B-\epsilon\,\Omega)\bigg(\partial_{\bm{\xi}}\frac{1}{n_0}\frac{\delta F}{\delta\varrho}\bigg)\cdot\mathbb{J}\bigg(\partial_{\bm{\xi}}\frac{1}{n_0}\frac{\delta G}{\delta\varrho}\bigg)\,n_0\,\varrho\,d\bm{\xi}\,d\bm{q}\nonumber\\
%%%
&-\epsilon\int \bigg(\left\langle\partial_{\bm{q}}\frac{1}{n_0}\frac{\delta F}{\delta\varrho}\right\rangle\cdot\left\langle\partial_{\bm{\xi}}\frac{1}{n_0}\frac{\delta G}{\delta\varrho}\right\rangle-\left\langle\partial_{\bm{q}}\frac{1}{n_0}\frac{\delta G}{\delta\varrho}\right\rangle\cdot\left\langle\partial_{\bm{\xi}}\frac{1}{n_0}\frac{\delta F}{\delta\varrho}\right\rangle\bigg)\,n\,d\bm{q}-\int (B-\epsilon\,\Omega)\, \left\langle\partial_{\bm{\xi}}\frac{1}{n_0}\frac{\delta F}{\delta\varrho}\right\rangle\cdot \mathbb{J} \left\langle\partial_{\bm{\xi}}\frac{1}{n_0}\frac{\delta G}{\delta\varrho}\right\rangle\,n_0\,d\bm{q}.\label{QNVP_cbracket}
\end{align}
Notice that the inverse Laplacian $\mathcal{G}$ does not appear in the bracket because of $\partial_{\bm{q}}\cdot (\delta G/\delta\bm{\pi}) = 0$. This bracket is polynomial in $(\bm{\pi},\varrho)$ and regular in $n_0$ away from $n_0 = 0$. It follows that $\Sigma$ is a Poisson-Dirac submanifold in $\mathfrak{s}^*$. 

\end{proof}

\end{document}
